% !TeX program = pdflatex
%==============================================================================
%  Trabalho Prático 01 --- Aulas 01 a 10 (sem as extras)
%
%  O roteiro de correção sai de "Trabalho pratico 01 - gabarito.tex", que apenas
%  liga a opção [gabarito] e inclui este arquivo --- fonte única, dois PDFs.
%  As caixas "solucao" e "criterio" só aparecem no roteiro; no enunciado o
%  comment.sty as descarta, e por isso \begin e \end delas ficam sozinhos na
%  linha.
%
%  Os números das soluções foram MEDIDOS no notebook
%  "Trabalho pratico 01 - resolucao.ipynb", nesta pasta, com scikit-learn 1.9.0,
%  e com as sementes que o enunciado fixa. Ao mexer numa semente, num dado ou
%  numa grade, rode o notebook de novo e confira as caixas.
%
%  Os dados são dois arquivos do UCI Machine Learning Repository, convertidos de
%  .xls para .csv SEM limpeza nenhuma (só os nomes das colunas do concreto foram
%  traduzidos, e a resposta do crédito foi renomeada para "default"). Ficam em
%  recursos/dados/concreto.csv e recursos/dados/credito.csv.
%==============================================================================
\providecommand{\gabaritoopt}{}
\documentclass[11pt,a4paper]{article}
\usepackage[\gabaritoopt]{../recursos/latex/estilo-avaliacao}

% Prazo e valor ficam aqui em cima, para serem trocados sem caçar no texto.
\newcommand{\prazo}{a definir pelo professor}
\newcommand{\valortotal}{10,0 pontos}

% \tarefa{código}{título}{valor}
\newcommand{\tarefa}[3]{%
  \bigskip\par\noindent
  {\large\bfseries\color{azulufrj}#1. #2}\hfill{\bfseries(#3)}\par
  \nopagebreak\vspace{0.3em}}

\newcommand{\cod}[1]{\texttt{#1}}

\begin{document}
\cabecalhotrabalho{01}{\valortotal}{\prazo}

%==============================================================================
\section*{O que é este trabalho}
%==============================================================================
Ao longo do curso você viu cada método isolado, quase sempre em dados simulados, em
que a resposta certa é conhecida. Neste trabalho a situação se inverte: dois
conjuntos de dados reais, públicos, de áreas diferentes, e nenhuma resposta
conhecida de antemão. O que se pede é o que se faz numa análise de verdade ---
entender os dados antes de modelar, montar um protocolo de avaliação honesto, comparar
métodos com as ferramentas certas e, no fim, \emph{chegar às suas próprias conclusões}.

O trabalho tem duas partes, uma de \textbf{regressão} e uma de
\textbf{classificação}, e cada tarefa corresponde a uma ou mais aulas do curso. A
maior parte do código de que você precisa já apareceu nas aulas práticas; o que é
novo aqui é juntar as peças e decidir, em cada passo, \emph{por que} fazer assim.

\begin{ideia}[o que vale nota]
Um número sem justificativa não será considerado. Quando uma tarefa pergunta ``por quê?'',
``o que isso significa?'' ou ``compare'', a resposta é um texto curto, em português,
numa célula de \emph{markdown} logo abaixo do código --- e é esse texto que recebe a
maior parte da nota. Um resultado errado com uma leitura honesta (``esperava $X$,
obtive $Y$, e desconfio de $Z$'') vale mais do que um resultado certo sem
explicação nenhuma.
\end{ideia}

%==============================================================================
\section{Regras de entrega}\label{sec:regras}
%==============================================================================
\begin{enumerate}[leftmargin=1.4cm]
  \item \textbf{Um único notebook} \cod{TP01\_<nome>.ipynb}, com as duas
        partes. Use o \cod{Trabalho pratico 01 - modelo.ipynb} como ponto de partida: ele já
        carrega os dados e tem uma seção por tarefa.
  \item O notebook deve \textbf{rodar do início ao fim sem erro} com
        \emph{Kernel $\to$ Restart \& Run All}, com os dois \cod{.csv} na mesma pasta
        dele, sem baixar nada da internet. Em um computador comum, a execução
        completa deve levar no máximo uns 20 minutos: se passar disso, reduza grades
        ou o número de árvores e diga que reduziu.
  \item \textbf{Bibliotecas}: \cod{numpy}, \cod{pandas}, \cod{matplotlib},
        \cod{scipy} e \cod{scikit-learn}, esta na versão 1.6 ou mais nova (o
        \cod{GroupKFold} de R2 usa o parâmetro \cod{shuffle}, que só existe a partir
        dela). Nada de outras bibliotecas de aprendizado de máquina.
  \item \textbf{Reprodutibilidade}: use as sementes do enunciado onde ele as fixa, e
        fixe você mesmo (\cod{random\_state} ou \cod{np.random.default\_rng}) toda
        outra aleatoriedade.
  \item \textbf{O conjunto de teste de cada parte deve ser usado uma única vez}, nas
        tarefas R8 e C7. Toda escolha anterior --- de modelo, de hiperparâmetro, de
        corte --- é feita por validação cruzada \emph{dentro do treino}. Usar o
        teste antes disso para escolher qualquer coisa invalida a tarefa
        correspondente.
  \item \textbf{Gráficos} com eixos rotulados e unidades. \textbf{Tabelas} com
        poucas casas decimais --- três ou quatro algarismos significativos bastam.
  \item O trabalho é \textbf{individual}. Pode (e deve) reaproveitar código das
        aulas práticas e das listas; o que se exige é que cada resposta escrita seja
        sua.
\end{enumerate}

\noindent
\textbf{Distribuição dos pontos.} Parte~I, 4,5 pontos; Parte~II, 4,5 pontos;
\textbf{qualidade do notebook}, 1,0 ponto --- roda do começo ao fim, está
organizado por tarefa, explica \emph{antes} de cada bloco de código o que ele vai
fazer, tem gráficos legíveis, etc.

\begin{criterio}
\textbf{Qualidade (1,0).} 0,3 roda sem erro do começo ao fim (se para num erro
pontual e o resto é coerente, 0,15); 0,3 organização por tarefa e
texto antes do código; 0,2 gráficos e tabelas legíveis; 0,2 teste tocado só em R8 e
C7. Quem escolheu modelo, hiperparâmetro ou corte olhando o teste perde esses 0,2
\emph{e} a nota da tarefa em que fez a escolha.
\end{criterio}

%==============================================================================
\section{Os dados}
%==============================================================================
Os dois conjuntos são do \emph{UCI Machine Learning Repository}
(\url{https://archive.ics.uci.edu}), um dos repositórios públicos de dados mais
usados em aprendizado de máquina; os dois também estão no Kaggle. Ambos são
distribuídos sob a licença CC~BY~4.0. Os arquivos que você recebe são os originais,
convertidos de planilha para \cod{.csv} \textbf{sem limpeza nenhuma}: só os nomes das
colunas do concreto foram traduzidos, e a resposta do crédito foi renomeada para
\cod{default}. Tudo o que houver de estranho nos dados está lá de propósito --- ou
melhor, está lá porque dados reais são assim.

\subsection{Parte I --- \texttt{concreto.csv}: resistência do concreto à compressão}

\noindent\textbf{Fonte.} I.-C. Yeh (1998), ``Modeling of strength of
high-performance concrete using artificial neural networks'', \emph{Cement and
Concrete Research} 28(12), 1797--1808. Repositório UCI, \emph{Concrete Compressive
Strength}, DOI \href{https://doi.org/10.24432/C5PK67}{10.24432/C5PK67}.

\medskip\noindent\textbf{O problema.} O concreto é o material mais usado na
engenharia civil, e a sua resistência à compressão é a propriedade que define para
que ele serve. Ela depende da \emph{receita} --- as quantidades de cimento, água,
aditivos e agregados em um metro cúbico de mistura, o que os engenheiros chamam de
\textbf{traço} --- e da \textbf{idade}: o concreto continua ganhando resistência por
meses depois de moldado. Medir a resistência exige moldar corpos de prova, esperar e
rompê-los numa prensa, o que é caro e lento --- a norma usual é esperar 28 dias. Um
bom preditor permitiria estimar a resistência de um traço \emph{novo} antes de
ensaiá-lo.

\medskip\noindent\textbf{As variáveis.} São $n=1030$ ensaios de laboratório, sem
valores faltantes:
\begin{center}\small
\renewcommand{\arraystretch}{1.2}
\begin{tabular}{@{}lll@{}}
\toprule
\textbf{coluna} & \textbf{significado} & \textbf{unidade} \\
\midrule
\cod{cimento}            & cimento Portland                               & kg/m$^3$ de mistura \\
\cod{escoria}            & escória de alto-forno (subproduto siderúrgico)  & kg/m$^3$ \\
\cod{cinza\_volante}     & cinza volante (subproduto de termelétricas)    & kg/m$^3$ \\
\cod{agua}               & água                                           & kg/m$^3$ \\
\cod{superplastificante} & aditivo que dá fluidez com menos água          & kg/m$^3$ \\
\cod{agregado\_graudo}   & brita                                          & kg/m$^3$ \\
\cod{agregado\_miudo}    & areia                                          & kg/m$^3$ \\
\cod{idade}              & idade do corpo de prova no ensaio              & dias (1 a 365) \\
\midrule
\cod{resistencia}        & \textbf{resposta}: resistência à compressão    & MPa \\
\bottomrule
\end{tabular}
\end{center}

\noindent
Um fato da engenharia que pode vir a ser útil: a \textbf{lei de Abrams} (1918) diz que, para
os mesmos materiais, a resistência cai com a \textbf{razão água/cimento} --- é a
regra mais antiga e mais citada da dosagem de concreto.

\subsection{Parte II --- \texttt{credito.csv}: inadimplência no cartão de crédito}

\noindent\textbf{Fonte.} I.-C. Yeh e C.-H. Lien (2009), ``The comparisons of data
mining techniques for the predictive accuracy of probability of default of credit
card clients'', \emph{Expert Systems with Applications} 36(2), 2473--2480.
Repositório UCI, \emph{Default of Credit Card Clients}, DOI
\href{https://doi.org/10.24432/C55S3H}{10.24432/C55S3H}.

\medskip\noindent\textbf{O problema.} São $n=30\,000$ clientes de um banco de
Taiwan, acompanhados de abril a setembro de 2005. Para cada um, sabe-se o limite de
crédito, alguns dados cadastrais e o histórico dos últimos seis meses --- se pagou
em dia ou com atraso, quanto veio na fatura e quanto pagou. A resposta é se o
cliente \textbf{ficou inadimplente no mês seguinte} (outubro de 2005). Os valores
estão em dólares de Taiwan (NT\$).
\begin{center}\small
\renewcommand{\arraystretch}{1.2}
\begin{tabular}{@{}lp{10.6cm}@{}}
\toprule
\textbf{coluna} & \textbf{significado} \\
\midrule
\cod{ID}          & número do cliente \\
\cod{LIMIT\_BAL}  & limite de crédito concedido (NT\$), individual mais o da família \\
\cod{SEX}         & 1 = masculino; 2 = feminino \\
\cod{EDUCATION}   & 1 = pós-graduação; 2 = universidade; 3 = ensino médio; 4 = outros \\
\cod{MARRIAGE}    & 1 = casado(a); 2 = solteiro(a); 3 = outros \\
\cod{AGE}         & idade (anos) \\
\cod{PAY\_0}, \cod{PAY\_2}, \dots, \cod{PAY\_6}
                  & situação do pagamento em setembro (\cod{PAY\_0}), agosto
                    (\cod{PAY\_2}), \dots, abril (\cod{PAY\_6}): $-1$ = pagou em dia;
                    $1$ = um mês de atraso; $2$ = dois meses; \dots; $9$ = nove meses
                    ou mais. (Não existe \cod{PAY\_1}: o nome é assim no original.) \\
\cod{BILL\_AMT1}, \dots, \cod{BILL\_AMT6}
                  & valor da fatura (NT\$) em setembro, agosto, \dots, abril \\
\cod{PAY\_AMT1}, \dots, \cod{PAY\_AMT6}
                  & valor pago (NT\$) em setembro, agosto, \dots, abril \\
\midrule
\cod{default}     & \textbf{resposta}: 1 se ficou inadimplente em outubro, 0 se não \\
\bottomrule
\end{tabular}
\end{center}

\noindent
A descrição acima é a da documentação oficial. Você vai encontrar nos dados códigos
que ela \emph{não} descreve; tratá-los é parte da Tarefa~C1.

%==============================================================================
\section{Parte I --- Regressão: a resistência do concreto \hfill (4,5 pontos)}
%==============================================================================

%------------------------------------------------------------------------------
\tarefa{R1}{Conhecendo os dados}{0,6 ponto}
%------------------------------------------------------------------------------
Antes de ajustar um modelo, vamos olhar para os dados.

\begin{itens}
  \item Quantas observações e covariáveis há? O problema é de regressão ou de
        classificação?
  \item Faça o histograma da resistência e os diagramas de dispersão da resistência
        contra a idade e contra a razão água/cimento. Qual das relações é mais
        visivelmente não linear? Em que idades os ensaios foram feitos?
  \item Procure \textbf{linhas duplicadas}, contando as que repetem uma linha anterior
        (é o que o \cod{duplicated()} do \cod{pandas} conta): quantas repetem uma linha
        anterior por inteiro, e quantas repetem só as oito covariáveis (estas incluem
        aquelas)? Uma
        resistência medida em laboratório, igual a outra até a oitava casa decimal,
        é um segundo ensaio ou outra coisa? Decida o que fazer com as linhas
        inteiramente repetidas, \textbf{justifique}, e siga com essa decisão no resto
        da Parte~I.
  \item Um grupo de linhas com as mesmas oito covariáveis é um conjunto de
        \textbf{réplicas} --- o mesmo experimento repetido ---, desde que cada linha seja
        um ensaio de verdade. Com esses grupos dá para estimar o erro irredutível
        $\sigma^2=\Var(Y\mid\X=\x)$ sem modelo nenhum, pela variância combinada dentro
        dos grupos,
        \[
          \widehat\sigma^2=\frac{\sum_g\sum_{i\in g}(Y_i-\bar Y_g)^2}{\sum_g(n_g-1)},
        \]
        em que $g$ percorre os grupos, $n_g$ é o tamanho do grupo e $\bar Y_g$ a sua
        média. Calcule $\widehat\sigma^2$ com todos os grupos de linhas com as mesmas
        oito covariáveis, \emph{antes} e \emph{depois} da decisão do item~(c), com os
        graus de liberdade $\sum_g(n_g-1)$ de cada caso, e
        um intervalo de 95\% para $\sigma^2$ supondo normalidade e variância
        constante. O que essa estimativa diz sobre o menor erro quadrático médio que
        \emph{qualquer} preditor poderia atingir?
  \item O mesmo \textbf{traço} (as sete quantidades de ingredientes) foi ensaiado em
        várias idades. Crie uma coluna que identifique o traço,
\begin{lstlisting}
componentes = ["cimento", "escoria", "cinza_volante", "agua",
               "superplastificante", "agregado_graudo", "agregado_miudo"]
concreto["traco"] = concreto.groupby(componentes).ngroup()
\end{lstlisting}
        e conte: quantos traços distintos há, e quantas linhas por traço? Por que
        isso significa que as linhas da tabela \emph{não} são observações
        independentes umas das outras?
\end{itens}

\begin{solucao}
\textbf{(a)} $n=1030$, $p=8$, todas quantitativas e sem faltantes. A resposta é
quantitativa --- faz sentido uma ``resistência média'' ---, então é regressão.

\textbf{(b)} A idade é a relação mais não linear: sobe depressa nas primeiras semanas e
quase estabiliza depois de 90 dias, o que sugere $\log(\text{idade})$. A razão
água/cimento tem relação decrescente e razoavelmente monótona (Abrams). Os ensaios
foram feitos em só 14 idades, de 1 a 365 dias; 28 dias tem 425 das 1030 linhas.

\textbf{(c)} 25 linhas são idênticas a uma anterior, e 34 repetem as oito
covariáveis. Uma medida contínua repetida até a oitava casa não é réplica: é registro
duplicado. O esperado é remover as 25, ficando com $n=1005$. (Manter e justificar ---
``não sei se são erros'' --- é aceitável, desde que a consequência para o item~(d) seja
discutida.)

\textbf{(d)} Com as duplicatas: 19 grupos, 34 graus de liberdade,
$\widehat\sigma^2=25{,}1$, IC~$(16{,}4;\ 43{,}1)$. Sem elas: 9 grupos, 9 graus de
liberdade, $\widehat\sigma^2=64{,}6$, IC~$(30{,}6;\ 215{,}5)$. Os zeros das duplicatas
puxavam a estimativa para baixo e davam uma falsa precisão; o que sobra é
praticamente inútil --- um intervalo que vai de 31 a 215. A leitura correta: o piso do
risco existe ($\E[\Var(Y\mid\X)]$), mas estes dados quase não o identificam;
e, com variância não constante, as poucas réplicas nem representariam todas as regiões
do espaço. (Na Tarefa~R8 os melhores modelos ficam com EQM de validação em torno de
35, compatível com o intervalo.)

\textbf{(e)} 428 traços para 1005 linhas; 246 traços aparecem uma vez e 77 aparecem
cinco vezes. Linhas do mesmo traço compartilham a receita --- e o que quer que tenha
acontecido no preparo daquela mistura ---, então não são independentes: o traço é a
unidade, a linha é uma medida repetida dela.
\end{solucao}

\begin{criterio}
0,1 por item, e 0,2 no (d). No (c), a nota vem da justificativa: remover sem dizer
por quê vale metade; contar todas as linhas envolvidas (36 e 53, em vez de 25 e 34),
dizendo que contou assim, também vale. No (d), exigir os dois cálculos e a leitura de que 9 graus de
liberdade não sustentam conclusão; o intervalo pode ser qualquer um corretamente
justificado.
\end{criterio}

%------------------------------------------------------------------------------
\tarefa{R2}{Protocolo de avaliação}{0,4 ponto}
%------------------------------------------------------------------------------
A pergunta que queremos responder é: \emph{qual a
resistência de um traço que ainda não foi ensaiado?}

\begin{itens}
  \item Separe o conjunto de teste \textbf{por traço}, exatamente assim:
\begin{lstlisting}
import sklearn.model_selection as skm
X = concreto[componentes + ["idade"]]
y = concreto["resistencia"]
separa = skm.GroupShuffleSplit(n_splits=1, test_size=0.25, random_state=2026)
i_tr, i_te = next(separa.split(X, y, groups=concreto["traco"]))
\end{lstlisting}
        O \cod{X} tem só as oito covariáveis originais: a coluna \cod{traco} identifica
        o grupo, não é covariável. Quantas linhas e quantos traços ficaram em cada lado?
        Confira que nenhum traço aparece nos dois.
  \item Por que separar por traço, e não por linha com \cod{train\_test\_split}? Que
        tipo de vazamento --- grave ou leve --- a separação por linha produziria,
        e para qual pergunta ela seria, ao contrário, a separação certa?
  \item Toda escolha de hiperparâmetro da Parte~I será feita por validação cruzada
        de 5 dobras \textbf{em grupo}, dentro do treino, passando \cod{groups=} os
        traços do treino:
\begin{lstlisting}
dobras = skm.GroupKFold(n_splits=5, shuffle=True, random_state=0)
\end{lstlisting}
        Reporte sempre o risco estimado como
        \textbf{média $\pm$ erro-padrão} das cinco dobras. O que esse erro-padrão
        mede, e o que ele \emph{não} mede?
\end{itens}

\begin{solucao}
\textbf{(a)} 321 traços e 756 linhas no treino, 107 traços e 249 linhas no teste (a
proporção de 25\% é de traços, não de linhas); nenhum traço em comum.

\textbf{(b)} Com a separação por linha, o mesmo traço, ensaiado em outra idade, ficaria
no treino e no teste --- ``a mesma unidade dos dois lados'', um vazamento
\textbf{grave}, que nenhum \emph{pipeline} conserta. O teste
mediria o acerto em traços já conhecidos. A separação por linha seria a certa para
outra pergunta: prever a resistência de um traço já ensaiado, numa idade que ainda não
foi.

\textbf{(c)} O erro-padrão mede a variabilidade entre as cinco dobras --- quanto a
estimativa dependeria de quais traços caíram em cada uma. Não mede a variabilidade da
amostra inteira (outro conjunto de 428 traços), e é otimista mesmo como medida da
primeira coisa, porque as dobras compartilham treino e as cinco estimativas não são
independentes.
\end{solucao}

\begin{criterio}
0,1 no (a); 0,2 no (b) --- 0,1 por nomear o vazamento, 0,1 por identificar a
pergunta para a qual a separação por linha serviria; 0,1 no (c). No (c), aceitar a
resposta sem a ressalva da dependência entre dobras, desde que diga o que o EP mede.
\end{criterio}

%------------------------------------------------------------------------------
\tarefa{R3}{Modelos lineares e regularização}{0,9 ponto}
%------------------------------------------------------------------------------

\begin{itens}
  \item Estime o risco do \textbf{preditor constante} (a média do treino,
        \cod{DummyRegressor}) e do \textbf{MQO} nas oito covariáveis originais. Some
        as sete colunas de ingredientes linha a linha: o que você observa? Calcule o
        número de condição de $\mathbb{X}^\top\mathbb{X}$ (com a coluna de uns) nas
        escalas original e padronizada. Com isso em mente, faz sentido ler o
        coeficiente da água como ``o efeito de aumentar a água em 1~kg/m$^3$
        mantendo as outras covariáveis fixas''?
  \item Troque a idade por $\log(\text{idade})$ e acrescente a razão água/cimento ---
        duas covariáveis construídas, nove ao todo --- e reestime o risco do MQO. Por
        que o modelo
        continua sendo uma regressão \emph{linear}? De onde veio o ganho, em termos
        de viés e variância?
  \item Sobre as covariáveis de (b), faça a expansão polinomial de grau~2
        (\cod{PolynomialFeatures}) e ajuste MQO, Ridge e Lasso, com a padronização
        \emph{dentro} do \emph{pipeline}. Escolha o $\lambda$ do Ridge e o do Lasso por
        validação cruzada numa grade logarítmica e mostre a curva do risco estimado contra
        $\lambda$, com barras de um erro-padrão. Aplique também a \textbf{regra de um
        erro-padrão}. Quantos coeficientes o Lasso zera em cada escolha? Faça o
        gráfico do caminho de regularização do Lasso.
  \item Por que foi preciso padronizar antes do Ridge e do Lasso, e não antes do
        MQO? O $\lambda$ escolhido para o Ridge e o escolhido para o Lasso podem ser
        comparados diretamente no \cod{scikit-learn}?
\end{itens}

\begin{solucao}
\textbf{(a)} Constante: $250{,}5\pm21{,}7$. MQO original: $105{,}4\pm5{,}6$. Os sete
ingredientes somam quase o mesmo valor em toda linha (média 2339, desvio
$63{,}8$~kg/m$^3$) --- é a massa de um metro cúbico de concreto. Com o intercepto, as
colunas ficam quase linearmente dependentes: número de condição $1{,}0\times10^{10}$
na escala original, $72$ padronizado (a padronização resolve a parte numérica, não a
colinearidade de fundo). ``Aumentar a água mantendo o resto fixo'' descreve uma
mistura que muda de volume --- que não existe; aumentar a água é, na prática,
diminuir outra coisa. O coeficiente é uma comparação entre misturas que diferem em
tudo, e não um efeito causal.

\textbf{(b)} $51{,}5\pm3{,}4$: o risco cai à metade. O modelo continua linear
\emph{nos parâmetros}; as covariáveis é que são funções dos dados. O ganho
é de \textbf{viés}: a classe de funções passou a conter formas (crescimento
logarítmico, efeito da razão) que a reta nas covariáveis originais não alcançava,
com só um parâmetro a mais.

\textbf{(c)} Com 54 covariáveis: MQO $53{,}7\pm7{,}9$; Ridge $43{,}8\pm4{,}3$ no
mínimo ($\alpha\approx4$) e $46{,}7\pm3{,}5$ pela regra 1-EP ($\alpha\approx63$); Lasso
$44{,}7\pm5{,}3$ no mínimo ($\alpha\approx0{,}012$, 37 de 54 coeficientes não nulos) e
$49{,}6\pm3{,}6$ pela regra 1-EP ($\alpha\approx0{,}18$, 23 não nulos). A expansão
sozinha piora o MQO por variância; a penalização recupera. Leituras esperadas: a regra
1-EP troca pouco risco por um modelo bem mais simples; o $\lambda$ que melhor prediz
não é o que melhor seleciona.

\textbf{(d)} A penalização soma $\beta_j^2$ ou $\abs{\beta_j}$, que dependem da
unidade da covariável; o MQO é equivariante (a predição não muda). E não: o
\cod{Ridge} minimiza $\norm{\bm y-\mathbb X\bbeta}^2+\alpha\norm{\bbeta}^2$ e o
\cod{Lasso}, $\frac1{2n}\norm{\bm y-\mathbb X\bbeta}^2+\alpha\norm{\bbeta}_1$ --- há um
fator $2n$ entre as escalas de $\alpha$.
\end{solucao}

\begin{criterio}
(a) 0,25 --- 0,05 pelos riscos, 0,1 pela soma e pelo número de condição, 0,1 pela
leitura do coeficiente. (b) 0,15. (c) 0,35 --- 0,1 pela validação sem vazamento
(padronização no \emph{pipeline}, dobras em grupo), 0,1 pelas curvas, 0,1 pela regra
1-EP e contagem de zeros, 0,05 pelo caminho. (d) 0,15. Grades que param na borda
(o $\lambda$ escolhido é o menor ou o maior da grade) sem comentário perdem 0,05. No
erro-padrão da regra 1-EP, o desvio das cinco dobras pode ter divisor 5 (o
\cod{std\_test\_score} do \cod{GridSearchCV}) ou 4 (o da solução): no Lasso, a regra
escolhe $\alpha\approx0{,}14$ num caso e $0{,}18$ no outro, e os dois valem.
\end{criterio}

%------------------------------------------------------------------------------
\tarefa{R4}{Validação cruzada, LOOCV e \emph{bootstrap}}{0,6 ponto}
%------------------------------------------------------------------------------
Use o MQO com as covariáveis de R3(b).

\begin{itens}
  \item Implemente à mão a validação cruzada em grupo, dobra a dobra, com as dobras
        de R2(c). As dobras têm o mesmo tamanho? Calcule o risco pela média
        \emph{simples} dos cinco EQM e pela média \emph{ponderada} pelo tamanho das
        dobras (a fórmula das notas de aula), e compare com o
        \cod{cross\_val\_score}.
  \item Calcule o LOOCV de duas formas: com $n$ ajustes, e pelo atalho da
        \emph{leverage} $h_{ii}$, com um único ajuste. Confira
        que coincidem. O LOOCV fica acima ou abaixo da validação em grupo? Por quê, e
        qual dos dois responde à pergunta de R2?
  \item Estime o erro-padrão de cada coeficiente de três formas: pela fórmula
        clássica, a raiz da diagonal de
        $\widehat\sigma^2(\mathbb{X}^\top\mathbb{X})^{-1}$, com
        $\widehat\sigma^2=\text{RSS}/(n-p-1)$; por \emph{bootstrap} das
        \textbf{linhas}; e por \emph{bootstrap} dos \textbf{traços} --- sorteando traços
        inteiros, com reposição, e levando todas as linhas de cada traço sorteado. Use
        $B=2000$ e o gerador \cod{np.random.default\_rng(2026)}. Compare as três
        colunas. Qual reamostragem imita o jeito como os dados foram coletados?
\end{itens}

\begin{solucao}
\textbf{(a)} Não: 160, 150, 153, 147 e 146 linhas, porque os traços têm tamanhos
diferentes. Média simples $51{,}526$ (igual ao \cod{cross\_val\_score}, que faz a
média não ponderada); ponderada $51{,}429$. A diferença é pequena, e existe porque
as duas fórmulas só coincidem quando as dobras têm o mesmo tamanho.

\textbf{(b)} Atalho e explícito coincidem até a sexta casa decimal: $49{,}360332$. O
LOOCV fica
\emph{abaixo} da validação em grupo ($51{,}53$): ao deixar uma linha de fora, as outras
idades do mesmo traço continuam no treino, e o modelo ``já conhece'' aquela receita.
O LOOCV responde a ``prever um ensaio novo de um traço conhecido''; a pergunta de R2 é
a da validação em grupo. Para o MQO, com 10 parâmetros, a diferença é pequena --- o
modelo quase não consegue decorar traços.

\textbf{(c)} Para $\log(\text{idade})$: $0{,}218$, $0{,}220$ e $0{,}304$; para a
razão água/cimento: $2{,}10$, $2{,}31$ e $3{,}45$. O \emph{bootstrap} das linhas
reproduz aproximadamente a fórmula clássica (e é um pouco maior, porque não supõe
variância constante); o dos traços dá erros-padrão de 1,4 a 1,6 vez os do
\emph{bootstrap} das linhas, e de 1,4 a 2,1 vezes os da fórmula clássica, conforme o
coeficiente. Os dados foram coletados por traço --- cada receita, ensaiada
em várias idades ---, e a reamostragem que imita isso é a dos traços. As outras duas
tratam as linhas como independentes e subestimam a incerteza.
\end{solucao}

\begin{criterio}
(a) 0,2; (b) 0,2 --- 0,1 pela igualdade, 0,1 pela comparação com a validação em
grupo; (c) 0,2 --- 0,1 pelos três cálculos, 0,1 pela leitura. Implementar o
\emph{bootstrap} dos traços errado (sorteando linhas dentro de traços, por exemplo)
vale metade do (c) se a leitura estiver certa.
\end{criterio}

%------------------------------------------------------------------------------
\tarefa{R5}{Métodos não paramétricos: KNN}{0,5 ponto}
%------------------------------------------------------------------------------

\begin{itens}
  \item Ajuste o KNN com as covariáveis padronizadas, escolhendo $k$ de 1 a 30 por
        validação cruzada, e mostre a curva do risco estimado contra $k$. Qual $k$
        foi escolhido? Ajuste também o KNN \emph{sem} padronizar, com o mesmo $k$, e
        explique a diferença (pense em quanto cada covariável pesa na distância, com e
        sem padronização).
\end{itens}

\begin{solucao}
\textbf{(a)} A validação escolhe $k=1$ ($57{,}3\pm4{,}1$), e a curva cresce com $k$
(a menos de uma oscilação desprezível em $k=14$) --- com poucos traços parecidos entre
si, o vizinho mais próximo é o único que ajuda. (O mesmo acontece com as covariáveis
construídas: $55{,}4$.) Sem padronizar, com o mesmo $k=1$, o risco sobe para
$74{,}1\pm3{,}2$. Sem padronização, cada covariável pesa na distância em proporção à
sua variância: o cimento responde por 29\% dela, e a água e o superplastificante, que
variam poucas dezenas de kg, por 1\% e 0,1\% --- a água, que a lei de Abrams põe no
centro, quase some da distância. Padronizadas, as oito pesam o mesmo. (Com $k=5$ a
diferença é bem menor, $88{,}0$ contra $85{,}4$, como mostra a tabela de R7.)
\end{solucao}

\begin{criterio}
0,25 pela curva e pela escolha de $k$; 0,25 pela explicação do que a padronização
muda na distância.
\end{criterio}

%------------------------------------------------------------------------------
\tarefa{R6}{Árvores, \emph{bagging} e florestas}{0,5 ponto}
%------------------------------------------------------------------------------

\begin{itens}
  \item Ajuste uma árvore de regressão cheia e calcule o seu erro de treino e o seu
        risco estimado. Depois, faça a poda por custo--complexidade, escolhendo o
        \cod{ccp\_alpha} por validação cruzada. Desenhe a árvore com profundidade~3
        (\cod{plot\_tree}) e interprete os três primeiros cortes.
  \item Ajuste o \emph{bagging} e florestas aleatórias com $m$ (\cod{max\_features})
        de 1 a 8. Mostre o risco estimado contra $m$. Qual é o padrão de
        \cod{max\_features} no \cod{RandomForestRegressor}, e o que isso quer dizer?
        Use o piso $\rho\,v$ da variância para explicar a forma da curva.
  \item Mostre o risco estimado da floresta escolhida contra o número de árvores
        $B$, de 1 a algumas centenas. Por que $B$ não é escolhido por validação
        cruzada como $m$?
\end{itens}

\begin{solucao}
\textbf{(a)} A árvore cheia tem 740 folhas e erro de treino $0{,}05$ --- não é zero
porque as réplicas, com as mesmas covariáveis e respostas diferentes, não podem ser
separadas ---, e risco estimado $56{,}4\pm2{,}8$. A poda escolhe um $\alpha$ pequeno
(320 folhas) e quase não ajuda: $55{,}9\pm2{,}5$. O primeiro corte é na idade, em 21
dias, e os dois seguintes são no cimento: em 386~kg/m$^3$ entre os corpos de prova de
até 21 dias e em 357,5~kg/m$^3$ entre os mais velhos. A idade vem primeiro porque é o
corte que mais separa as resistências --- um concreto jovem ainda não chegou à
resistência que terá ---, e nos dois ramos mais cimento dá mais resistência.

\textbf{(b)} De $47{,}3$ ($m=1$) o risco cai até $35{,}5$ ($m=5$) e fica praticamente
plano até o \emph{bagging} ($m=8$: $36{,}3$); todas as diferenças de $m=3$ em diante
estão dentro de um erro-padrão (cerca de $4{,}3$). O padrão do
\cod{RandomForestRegressor} é \cod{max\_features=1.0}, isto é, todas as covariáveis:
é \emph{bagging}, e não floresta. Com $m=1$ as árvores ficam descorrelacionadas mas
fracas demais (o $v$ e o viés sobem mais do que o $\rho$ cai); com $p=8$ covariáveis,
das quais poucas dominam, há pouco a descorrelacionar.

\textbf{(c)} De $87{,}9$ com uma árvore o risco desce a $38{,}6$ com 10, $35{,}9$ com
50 e estaciona em $35{,}2$--$35{,}5$ de 100 a 400. Aumentar $B$ só reduz o termo
$(1-\rho)v/B$, e não há superajuste em $B$: basta um $B$ grande o bastante.
\end{solucao}

\begin{criterio}
(a) 0,2 --- 0,1 pelos números e pela poda, 0,1 pela interpretação da árvore;
(b) 0,2 --- inclui dizer que o padrão é \emph{bagging}; (c) 0,1.
\end{criterio}

%------------------------------------------------------------------------------
\tarefa{R7}{Sensibilidade à escala}{0,5 ponto}
%------------------------------------------------------------------------------

\begin{itens}
  \item Suponha que os mesmos dados chegassem com a idade em \textbf{horas}, e não em
        dias: crie uma cópia das covariáveis de treino com a coluna \cod{idade}
        multiplicada por 24, sem mexer em mais nada. Estime o risco dos seis modelos
        abaixo, todos nas oito covariáveis originais, \emph{duas vezes}: com a idade em
        dias e com a idade em horas. Use a validação cruzada em grupo de R2(c) e
        mantenha fixos os hiperparâmetros --- nada é reescolhido na segunda estimativa.
        \begin{itemize}[nosep]
          \item o MQO;
          \item o Lasso com \cod{alpha=1}, com as covariáveis padronizadas dentro do
                \emph{pipeline};
          \item o KNN com $k=5$, sem padronizar;
          \item o mesmo KNN, com as covariáveis padronizadas dentro do \emph{pipeline};
          \item a árvore podada de R6(a), com o \cod{ccp\_alpha} escolhido lá;
          \item a floresta de R6(b), com o $m$ escolhido lá.
        \end{itemize}
        Monte uma tabela com uma linha por modelo e uma coluna para cada unidade. Quais
        riscos mudaram com a troca de unidade, e quais ficaram iguais? Explique, modelo a
        modelo, por que a unidade da idade afeta ou não o resultado.
\end{itens}

\begin{solucao}
\textbf{(a)} Só o KNN sem padronizar muda: de $88{,}0$ para $63{,}4$. MQO
($105{,}4$), Lasso com \cod{alpha=1} ($119{,}3$), KNN padronizado ($85{,}4$), árvore e floresta
ficam praticamente idênticos (o KNN padronizado oscila na terceira casa decimal, de
$85{,}352$ para $85{,}351$). O MQO é equivariante à escala; o Lasso e o KNN
padronizados absorvem a mudança no \cod{StandardScaler}; árvore e floresta comparam
valores dentro de cada coluna. O KNN cru é o único que enxerga a unidade --- e aqui
melhora por acaso, porque em horas a idade (a covariável mais informativa) passa a
dominar a distância. Quem decidiu foi a unidade de medida, não o analista.
\end{solucao}

\begin{criterio}
0,25 pela tabela dos doze riscos; 0,25 pela explicação, modelo a modelo, de quem muda e
por quê. Não é preciso usar nenhuma classificação dos métodos: basta dizer, de cada um,
por que a unidade não o afeta (ou afeta).
\end{criterio}

%------------------------------------------------------------------------------
\tarefa{R8}{Escolha final e teste}{0,5 ponto}
%------------------------------------------------------------------------------

\begin{itens}
  \item Monte uma tabela com o risco estimado ($\pm$ erro-padrão) de pelo menos um
        modelo de cada família: constante, linear regularizado, KNN, árvore e
        floresta. Escolha \textbf{um} modelo final e
        justifique --- a regra de um erro-padrão pode entrar no argumento.
  \item Reajuste o modelo escolhido em \emph{todo} o treino e meça, \textbf{uma
        vez}, o EQM no teste, com um intervalo de 95\% para o risco. O EQM de
        teste ficou acima ou abaixo da estimativa da validação cruzada?
  \item Só agora, calcule o EQM de teste dos outros finalistas --- para leitura, sem
        trocar a escolha. Os intervalos se sobrepõem? Algum modelo foi muito pior no
        teste do que a validação previa? Investigue os maiores erros dele: as
        covariáveis dessas observações estão dentro do intervalo visto no treino?
  \item Em um parágrafo, tente explicar para alguém que não fez este curso: o que o modelo
        consegue prever, com que erro típico (em MPa), e em que situação você não
        confiaria nele.
\end{itens}

\begin{solucao}
\textbf{(a)} Na validação: floresta ($m=5$) $35{,}5\pm4{,}5$, \emph{bagging}
$36{,}3\pm4{,}5$, Ridge grau~2 $43{,}8$, Lasso $44{,}7$, MQO com covariáveis
construídas $51{,}5$, KNN e árvore $55$--$57$, constante $250$. Floresta e
\emph{bagging} empatam dentro de um erro-padrão, e nenhum modelo mais simples fica a
menos de um erro-padrão do mínimo (a Ridge de grau~2, com $43{,}8$, passa do limite
$35{,}5+4{,}5=40{,}0$): a regra de um erro-padrão não troca a escolha. A resolução
fica com a floresta, o menor risco estimado; o \emph{bagging}, bem justificado, também
é aceitável.

\textbf{(b)} Floresta no teste: EQM $45{,}8$, IC~$(31{,}9;\ 59{,}7)$. O EQM ficou
\emph{acima} da estimativa da validação cruzada ($35{,}5$), mas o intervalo contém o
$35{,}5$: o teste não desmente a validação, e a diferença cabe na variabilidade de
sortear 107 traços. Quem quiser ir além acha a causa na variância da resistência,
$310{,}6$ no teste contra $248{,}2$ no treino: uma razão de $1{,}25$, quase a mesma com
que a floresta piora ($1{,}29$). O teste é mais difícil, e não há sinal de superajuste.

\textbf{(c)} No teste, com a meia largura do intervalo de 95\% depois do $\pm$:
\emph{bagging} $50{,}0\pm14{,}8$, Lasso 1-EP $50{,}0\pm11{,}8$, MQO
construído $56{,}4\pm12{,}3$, KNN $65{,}4\pm15{,}6$, árvore $92{,}7\pm25{,}0$: os
intervalos dos bons modelos se sobrepõem, e o teste não os distingue. Quem foi muito
pior do que a validação previa é a árvore podada: $55{,}9$ na validação e $92{,}7$ no
teste, uma razão de $1{,}66$, contra $1{,}01$ a $1{,}38$ dos outros finalistas. Os seis
maiores erros dela, que somam 24\% do EQM, são de dois traços do teste com a mesma
receita-base (água 228, os mesmos agregados, cimento e escória somando
380~kg/m$^3$), ensaiados de 180 a 365 dias --- e as covariáveis estão todas
\emph{dentro} do intervalo do treino: não é extrapolação. A árvore põe os nove
ensaios desses dois traços numa folha com só três linhas do treino, de outros traços,
todas de 28 dias e com cinza volante, de média $23{,}5$~MPa, quando eles resistem de
48 a 56~MPa. Os dois traços do treino com a mesma receita-base ficaram do outro lado
de um corte na escória ($\le97{,}5$), em folhas que preveem de 41 a 52~MPa. É a
instabilidade da árvore: um único corte, numa região com poucos traços, decide a
vizinhança. A única extrapolação do teste é um traço (5 linhas) com agregado miúdo
($992{,}6$) acima do máximo do treino ($945$), e a floresta não sofre com ele: sem essas
5 linhas, o EQM dela é $46{,}5$, praticamente o mesmo.

\textbf{(d)} Algo como: ``o modelo prevê a resistência de um concreto com uma receita
que ainda não foi ensaiada, e erra tipicamente uns 7~MPa, para mais ou para menos (a
raiz do EQM de teste); não confiaria nele para receitas fora da faixa das que foram
ensaiadas --- por exemplo, com mais areia do que qualquer mistura dos dados ---, nem
para idades fora de 1 a 365 dias.''
\end{solucao}

\begin{criterio}
(a) 0,15; (b) 0,15 --- o EQM com o intervalo, e a comparação com a validação lida
à luz do intervalo, não só ``ficou acima''; comparar as variâncias não é exigido, mas
vale como justificativa; (c) 0,1 ---
exigir que o aluno identifique o modelo que mais piorou da validação para o teste e
\emph{investigue} os maiores erros dele; concluir, com os números, que as covariáveis
estão dentro do intervalo do treino e que o problema é outro vale a nota cheia;
(d) 0,1 --- erro em MPa, não em MPa$^2$.
\end{criterio}

%==============================================================================
\section{Parte II --- Classificação: inadimplência no cartão de crédito \hfill (4,5 pontos)}
%==============================================================================

%------------------------------------------------------------------------------
\tarefa{C1}{Conhecendo os dados}{0,4 ponto}
%------------------------------------------------------------------------------

\begin{itens}
  \item Qual a proporção de inadimplentes? Qual a acurácia do classificador que diz
        ``ninguém fica inadimplente''? O que isso antecipa sobre a acurácia como
        métrica neste problema?
  \item Tabele os valores de \cod{SEX}, \cod{EDUCATION}, \cod{MARRIAGE} e
        \cod{PAY\_0}. Que códigos aparecem sem estar na documentação? Decida o que
        fazer com eles e justifique. Quais dessas colunas são nominais, e o que isso
        exige antes de usá-las num modelo que soma covariáveis? E o \cod{ID}?
  \item Os \cod{PAY\_*} misturam atrasos (1 a 9 meses, uma escala ordinal) com
        códigos ($-2$, $-1$, $0$) que não são ``quantidades de atraso''. Que problema
        isso cria para um modelo linear que usa \cod{PAY\_0} como número? Proponha
        uma alternativa.
  \item Calcule a matriz de correlação das seis colunas \cod{BILL\_AMT}. Para qual
        dos classificadores generativos isso é uma má notícia, e por quê?
  \item \cod{SEX}, \cod{AGE} e \cod{MARRIAGE} são covariáveis disponíveis. Em uma
        linha: vê algum problema em usá-las para decidir quem recebe crédito?
        (Não é preciso retirá-las do modelo; é preciso pensar nisso.)
\end{itens}

\begin{solucao}
\textbf{(a)} $22{,}12\%$; a acurácia de ``ninguém'' é $77{,}9\%$. Qualquer modelo tem
de ser comparado com esse piso, e mesmo um bom modelo vai ganhar poucos pontos de
acurácia (C4).

\textbf{(b)} \cod{EDUCATION} tem 0, 5 e 6 (345 clientes) e \cod{MARRIAGE} tem 0 (54),
não documentados; \cod{PAY\_0} tem $-2$ e $0$, também não documentados. Juntar os
códigos desconhecidos em ``outros'' é a escolha mais simples e defensável; tratá-los
como faltantes também. \cod{SEX} e \cod{MARRIAGE} são nominais; \cod{EDUCATION} tem
ordem nos códigos 1 a 3, mas o ``outros'' a quebra, e tratá-la também como nominal é o
mais simples (tratá-la como ordinal, dizendo o que se fez com o ``outros'', também
vale). As nominais precisam de \emph{dummies} (\cod{OneHotEncoder} num
\cod{ColumnTransformer}); usá-las como número impõe uma ordem e distâncias que não
existem. O \cod{ID} é
identificador e sai: para um cliente novo ele não quer dizer nada. (A inadimplência
varia entre os decis do \cod{ID} mais do que o acaso explicaria --- qui-quadrado
$51{,}7$ com 9 graus de liberdade ---, mas sem tendência. O \cod{ID} é a posição no
arquivo, e o que ele captura é o jeito como o arquivo foi montado, não uma propriedade
do cliente.)

\textbf{(c)} Como número, \cod{PAY\_0} supõe que a diferença entre $-2$ e $-1$ é a mesma
entre $1$ e $2$ meses de atraso, e que o efeito no \emph{log-odds} é linear na
escala. Alternativa: \emph{dummies} por código (medido em C3: a AUC da logística vai
de $0{,}719$ para $0{,}765$), ou uma indicadora ``em atraso'' mais o número de meses.

\textbf{(d)} Correlações de $0{,}80$ a $0{,}95$. É má notícia para o \textbf{Bayes
ingênuo}: ele trata seis colunas que dizem quase a mesma coisa como seis testemunhas
independentes, e as probabilidades saem extremas (medido em C4).

\textbf{(e)} Sim: decisão de crédito com base em sexo ou estado civil pode ser
discriminatória, e é restrita por lei em muitos países --- e retirar a coluna não
basta se outras covariáveis a reproduzirem. Qualquer resposta que reconheça a questão
vale a nota.
\end{solucao}

\begin{criterio}
0,1 em (a), (b) e (c); 0,05 em (d) e (e).
\end{criterio}

%------------------------------------------------------------------------------
\tarefa{C2}{Protocolo e custos}{0,4 ponto}
%------------------------------------------------------------------------------

\begin{itens}
  \item Separe o teste de forma estratificada, exatamente assim, e use estas
        dobras em toda validação cruzada da Parte~II (os nomes com \cod{c} evitam
        sobrescrever as variáveis da Parte~I):
\begin{lstlisting}
Xc_tr, Xc_te, yc_tr, yc_te = skm.train_test_split(
    Xc, yc, test_size=0.25, stratify=yc, random_state=2026)
dobras_c = skm.StratifiedKFold(5, shuffle=True, random_state=0)
\end{lstlisting}
        Por que estratificar aqui, se na Parte~I a separação foi por grupo? Há alguma
        unidade repetida nestes dados?
  \item O banco vai usar o modelo para decidir \textbf{quem contatar} com uma
        proposta de renegociação antes do vencimento. Contatar um cliente que pagaria
        de qualquer jeito custa $c_{FP}=1$ (tempo da equipe, desconto concedido sem
        necessidade); deixar de contatar um cliente que vai ficar inadimplente custa
        $c_{FN}=5$. Deduza o corte ótimo
        para a probabilidade \emph{verdadeira} $\Prob(Y=1\mid\x)$.
  \item Calcule, no treino, o custo médio por cliente das duas regras triviais,
        ``não contatar ninguém'' e ``contatar todos''. Qual delas é a referência a
        bater?
\end{itens}

\begin{solucao}
\textbf{(a)} A estratificação mantém a prevalência de $22{,}12\%$ nos dois lados e nas
dobras, o que estabiliza as métricas com classe desbalanceada. Aqui a unidade é o
cliente e cada um aparece uma vez (35 linhas são idênticas em tudo menos no
\cod{ID}, provavelmente clientes sem movimento; não é o caso de agrupar).

\textbf{(b)} Prever 1 custa $c_{FP}(1-\eta)$ em esperança; prever 0 custa
$c_{FN}\eta$. Prever 1 quando $c_{FP}(1-\eta)<c_{FN}\eta$, isto é,
$\eta>c_{FP}/(c_{FP}+c_{FN})=1/6\approx0{,}167$.

\textbf{(c)} Ninguém: $5\times0{,}2212=1{,}106$ por cliente. Todos:
$1\times0{,}7788=0{,}779$. A referência é ``contatar todos'' --- e um modelo só vale a
pena se ficar abaixo de $0{,}779$.
\end{solucao}

\begin{criterio}
(a) 0,1; (b) 0,2 --- a dedução, não só o número; (c) 0,1.
\end{criterio}

%------------------------------------------------------------------------------
\tarefa{C3}{Regressão logística e modelos generativos}{0,8 ponto}
%------------------------------------------------------------------------------
Em toda esta parte, o pré-processamento (padronização das
numéricas, \emph{dummies} das nominais) mora dentro do \emph{pipeline}.

\begin{itens}
  \item Ajuste a regressão logística com \cod{PAY\_*} como número, escolhendo o
        \cod{C} por validação cruzada com \cod{scoring="neg\_log\_loss"}. Reporte a
        razão de chances $e^{\beta_j}$ de \cod{PAY\_0}, de \cod{LIMIT\_BAL} e de
        \cod{BILL\_AMT1} e \cod{BILL\_AMT2}, e interprete as duas primeiras --- atenção
        à escala em que os coeficientes estão. O que acontece com os sinais de
        \cod{BILL\_AMT1} e \cod{BILL\_AMT2}, e por que eles não devem ser lidos
        isoladamente?
  \item Ajuste a mesma logística com os \cod{PAY\_*} em \emph{dummies} e compare as
        duas pela AUC, fora da dobra (\cod{cross\_val\_predict}).
  \item Ajuste o Bayes ingênuo gaussiano, o LDA e o QDA. O QDA com o
        \cod{reg\_param} padrão ajusta? Leia a mensagem de erro, explique a causa, e
        escolha um \cod{reg\_param} por validação cruzada. Quantos parâmetros de
        covariância o LDA e o QDA estimam no seu $p$ (depois das \emph{dummies})?
  \item Faça o histograma de \cod{LIMIT\_BAL} e de \cod{BILL\_AMT1} em cada classe.
        Quais suposições dos modelos generativos são claramente falsas aqui? Isso os
        impede de classificar bem? (Guarde a pergunta; C4 responde.)
\end{itens}

\begin{solucao}
\textbf{(a)} $C\approx0{,}32$. Com as numéricas padronizadas, $e^{\beta_j}$ é o fator
na chance por \textbf{um desvio-padrão} a mais: \cod{PAY\_0}, $1{,}85$ por DP
($\approx1{,}12$ ``unidades'' de código) --- quem atrasa mais tem uma chance quase duas
vezes maior; \cod{LIMIT\_BAL}, $0{,}92$ por DP ($\approx$ NT\$130 mil) --- limite maior,
chance um pouco menor, com o resto fixo. \cod{BILL\_AMT1} tem razão $0{,}64$ e
\cod{BILL\_AMT2}, $1{,}25$: sinais opostos em duas colunas com correlação $0{,}95$. É o
efeito da colinearidade --- o modelo contrasta as duas faturas, e ``aumentar uma
mantendo a outra fixa'' é uma situação quase inexistente nos dados. Lê-se o par, não
cada um.

\textbf{(b)} AUC fora da dobra: $0{,}719$ com \cod{PAY\_*} numérico e $0{,}765$ com
\emph{dummies} ($C=0{,}1$). A forma funcional estava errada, e a correção é de
\emph{covariáveis}, não de método.

\textbf{(c)} Bayes ingênuo: AUC $0{,}726$. LDA: $0{,}713$. O QDA sem regularização
levanta \cod{LinAlgError}: ``The covariance matrix of class 0 is not full rank''. As
\emph{dummies} de uma mesma variável somam (quase) uma constante e, dentro de uma
classe, algumas colunas ficam constantes ou colineares; a covariância da classe é
singular. Escolhido pelo \emph{log-loss}, o mesmo critério do \cod{C} da logística,
numa grade até $0{,}99$, o \cod{reg\_param} é $0{,}8$: um mínimo no interior da grade
(\emph{log-loss} $0{,}604$), com AUC $0{,}714$. Pela AUC a escolha seria $0{,}3$
($0{,}733$); qualquer um dos dois critérios vale, desde que o aluno diga qual usou.
Com $p=28$ depois das \emph{dummies}:
$p(p+1)/2=406$ parâmetros de covariância no LDA e $812$ no QDA.

\textbf{(d)} As duas são muito assimétricas, positivas, com massa em zero (faturas
nulas) e caudas longas; a normalidade dentro de cada classe é falsa, e a
independência do Bayes ingênuo também (C1). Mesmo assim os três ordenam razoavelmente
(AUC 0,71--0,73): classificar bem exige menos do que estimar bem as densidades. O
preço aparece nas probabilidades, em C4.
\end{solucao}

\begin{criterio}
(a) 0,25 --- 0,1 pela interpretação por desvio-padrão (interpretar como ``por
unidade'' com dados padronizados perde esses 0,1), 0,15 pela leitura do par de
faturas; (b) 0,15; (c) 0,25 --- 0,1 pela causa da recusa do QDA, 0,1 pelo
\cod{reg\_param} escolhido por validação, 0,05 pela conta; (d) 0,15.
\end{criterio}

%------------------------------------------------------------------------------
\tarefa{C4}{Métricas, custos e calibração}{1,0 ponto}
%------------------------------------------------------------------------------
Toda esta tarefa usa probabilidades \textbf{fora da dobra}
no treino, e o teste continua guardado:
\begin{lstlisting}
p = skm.cross_val_predict(modelo, Xc_tr, yc_tr, cv=dobras_c,
                          method="predict_proba")[:, 1]
\end{lstlisting}

\begin{itens}
  \item Para a logística com \cod{PAY\_*} em \emph{dummies}, no corte $0{,}5$: a matriz
        de confusão, a acurácia, a sensibilidade, a especificidade, a precisão, o VPN
        e o $F_1$. Compare a acurácia com a de C1(a). Qual erro o modelo comete mais?
  \item Sobreponha as curvas ROC e as curvas precisão--\emph{recall} da logística de
        (a), do Bayes ingênuo e do LDA, com a AUC e a precisão média de cada um.
  \item Compare a \textbf{calibração}: o score de Brier e a curva de calibração da
        logística e do Bayes ingênuo. Que fração das probabilidades do Bayes ingênuo
        passa de $0{,}99$? Calibre-o com \cod{CalibratedClassifierCV}
        (\cod{method="isotonic"}) e compare de novo. O que mudou na AUC, e o que mudou
        no Brier? Por quê?
  \item Para cada modelo até aqui, calcule o custo médio por cliente no corte
        $1/6$ de C2(b) e o menor custo possível varrendo o corte de $0{,}02$ a
        $0{,}95$. Em quais modelos o melhor corte fica perto de $1/6$, e em quais fica
        longe? Relacione com o item~(c).
  \item Reajuste a logística com \cod{class\_weight="balanced"}. O que acontece com o
        Brier e com o custo no corte $1/6$? Mostre que, se a reponderação
        multiplica a chance estimada por $n_0/n_1$, o corte que corresponde a $1/6$ no
        modelo reponderado é
        $\big[\tfrac15\cdot\tfrac{n_0}{n_1}\big]\big/\big[1+\tfrac15\cdot\tfrac{n_0}{n_1}\big]$,
        calcule-o e confira com o melhor corte que você encontrou.
\end{itens}

\begin{solucao}
\textbf{(a)} VN 16725, FP 798, FN 3248, VP 1729: acurácia $0{,}820$ (contra $0{,}779$ do
trivial --- quatro pontos), sensibilidade $0{,}347$, especificidade $0{,}955$,
precisão $0{,}684$, VPN $0{,}837$, $F_1$ $0{,}461$. O erro dominante é o falso
negativo: dois em cada três inadimplentes passam sem ser apontados no corte $0{,}5$ ---
justamente o erro que custa 5.

\textbf{(b)} AUC: logística $0{,}765$, Bayes ingênuo $0{,}726$, LDA $0{,}713$; precisão
média $0{,}536$, $0{,}479$ e $0{,}495$ (referência: a prevalência, $0{,}221$). A ROC da
logística domina as outras em quase toda a extensão.

\textbf{(c)} Brier: logística $0{,}137$, Bayes ingênuo $0{,}335$ --- quase duas vezes e
meia pior, com AUC comparável. $14{,}5\%$ das probabilidades do ingênuo passam de
$0{,}99$ (e $11\%$ ficam abaixo de $0{,}01$); a curva de calibração dele é quase
horizontal. Calibrado: Brier $0{,}145$, AUC $0{,}729$ --- a AUC quase não se mexe,
porque a calibração isotônica é monótona e preserva a ordem; o Brier cai porque os
números passam a significar frequências. Classificar (ordenar) bem $\neq$ estimar bem
probabilidades, e a causa é a correlação das faturas (C1d).

\textbf{(d)} No corte $1/6$ e no melhor corte: logística com \emph{dummies} $0{,}580$
e $0{,}579$ (corte $0{,}16$); Bayes ingênuo calibrado $0{,}618$ e $0{,}617$ ($0{,}17$);
logística com \cod{PAY\_*} numérico $0{,}691$ e $0{,}635$ ($0{,}23$); LDA $0{,}690$ e
$0{,}638$ ($0{,}21$); QDA $0{,}690$ e $0{,}647$ ($0{,}31$); Bayes ingênuo $0{,}701$ e
$0{,}622$ ($0{,}77$). O corte teórico só é o ótimo quando as probabilidades são
honestas: os modelos calibrados têm melhor corte em $1/6$; os descalibrados, longe. E
a logística com \cod{PAY\_*} numérico mostra que especificação errada também
descalibra.

\textbf{(e)} Brier $0{,}186$ (piora) e custo em $1/6$ de $0{,}755$ --- quase o de
``contatar todos''. Com $n_0/n_1=3{,}52$, a chance $1/5$ vira $0{,}704$ e o corte
equivalente é $0{,}413$; o melhor corte medido é $0{,}40$, com custo $0{,}579$ ---
praticamente o da logística sem peso no corte $1/6$, $0{,}580$. Reponderar é mover o corte por outro
caminho, e ao preço de estragar as probabilidades.
\end{solucao}

\begin{criterio}
(a) 0,15; (b) 0,15; (c) 0,3 --- 0,1 pelas medidas, 0,1 pela calibração do ingênuo,
0,1 pela explicação de por que a AUC não muda; (d) 0,25 --- a relação entre
calibração e corte é o centro da tarefa; (e) 0,15 --- 0,1 pela conta do corte
equivalente. Custos calculados no teste, em vez de fora da dobra, zeram o (d).
\end{criterio}

%------------------------------------------------------------------------------
\tarefa{C5}{SVM}{0,5 ponto}
%------------------------------------------------------------------------------
O custo de treinar o \cod{SVC} cresce entre $n^2$ e $n^3$. Faça esta
tarefa numa subamostra estratificada de 5000 clientes do treino:
\begin{lstlisting}
i_sub = skm.train_test_split(np.arange(len(yc_tr)), train_size=5000,
                             stratify=yc_tr, random_state=2026)[0]
\end{lstlisting}

\begin{itens}
  \item Ajuste a SVM com \emph{kernel} RBF, escolhendo \cod{C} e \cod{gamma} numa
        grade bidimensional por validação cruzada (\cod{scoring="roc\_auc"}), e a SVM
        linear. Mostre a tabela da grade. O ótimo ficou no interior dela?
  \item Compare com a logística ajustada na \emph{mesma} subamostra. Como a SVM não
        estima probabilidades, de onde vem a AUC dela?
  \item Conte os vetores de suporte para alguns valores de \cod{C}. No
        \cod{scikit-learn}, o que se espera que aconteça com eles quando \cod{C}
        cresce? Aconteceu? Explique o que viu pensando na sobreposição das classes.
  \item Seria razoável usar o corte $1/6$ de C2 nas ``probabilidades'' de
        \cod{SVC(probability=True)}? Por quê?
\end{itens}

\begin{solucao}
\textbf{(a)} Com \cod{C} em $\{0{,}1;1;10;100\}$ e \cod{gamma} em
$\{10^{-4},\dots,10^{-2}\}$, o ótimo é interior, perto de \cod{C}$=10$,
\cod{gamma}$=3\times10^{-4}$, AUC $\approx0{,}748$; o linear fica em $\approx0{,}74$.
Uma grade que ponha o ótimo na borda precisa ser estendida.

\textbf{(b)} A logística na mesma subamostra: AUC $0{,}753$ --- a SVM não ganha. A AUC
só precisa de um \emph{escore} que ordene, e a \cod{decision\_function} (a distância
com sinal à fronteira) serve.

\textbf{(c)} No \cod{scikit-learn} \cod{C} é o peso da penalidade: \cod{C} maior
tolera menos violações, a margem estreita, e espera-se \emph{menos} vetores de
suporte. Aqui eles ficam entre 2235 e 2336 de 5000 para qualquer \cod{C}, e chegam a
subir de \cod{C}$=0{,}1$ para \cod{C}$=10$. As classes se sobrepõem muito (o erro de Bayes aqui não é pequeno): todo ponto
mal classificado é vetor de suporte, e os pontos dentro da margem também; com tanta
sobreposição, quase metade da amostra está sempre nessa situação.

\textbf{(d)} Não sem verificar: as ``probabilidades'' do \cod{SVC} vêm de um ajuste de
Platt \emph{a posteriori}, feito por validação cruzada interna, e não do método. Se
estiverem calibradas o corte vale; é preciso conferir (Brier, curva de calibração)
antes de usá-lo.
\end{solucao}

\begin{criterio}
(a) 0,15; (b) 0,1; (c) 0,15 --- a nota é da explicação, e não de o número de vetores
ter caído; (d) 0,1.
\end{criterio}

%------------------------------------------------------------------------------
\tarefa{C6}{KNN, árvores e florestas}{0,6 ponto}
%------------------------------------------------------------------------------

\begin{itens}
  \item Ajuste o KNN (padronizado, dentro do \emph{pipeline}), escolhendo $k$ ímpar por
        validação cruzada pela AUC. Por que o $k$ escolhido é tão maior do que o da
        Parte~I?
  \item Ajuste árvores de classificação com Gini e com entropia, de profundidade 3:
        os cortes coincidem? Depois, faça a poda de uma árvore cheia por
        custo--complexidade, escolhendo o \cod{ccp\_alpha} por validação cruzada com
        \cod{scoring="neg\_log\_loss"}. Desenhe a árvore de Gini de profundidade~3 e
        interprete o primeiro corte.
  \item Ajuste florestas com $m=\sqrt p$ (o padrão do \cod{RandomForestClassifier}),
        $m\approx p/3$ e $m=p$ (\emph{bagging}), com \cod{min\_samples\_leaf=5}, e uma
        com $m=\sqrt p$ e folhas cheias, isto é, com o \cod{min\_samples\_leaf} padrão.
        Compare AUC e Brier. Opcional: o \cod{AdaBoostClassifier} e o custo dele no
        corte $1/6$.
\end{itens}

\begin{solucao}
\textbf{(a)} A AUC sobe de $0{,}613$ ($k=1$) a $0{,}745$ ($k=101$) e fica plana
($0{,}747$--$0{,}748$) de $k=151$ a $401$; a busca escolhe $k=301$, mas qualquer $k$
nesse platô é equivalente. A escolha é
grande porque o sinal é fraco e ruidoso (as classes se sobrepõem, $n$ é grande), e
estimar uma proporção local exige muitos vizinhos; na Parte~I o sinal era forte e
$n$, pequeno.

\textbf{(b)} Os cortes das duas árvores de profundidade 3 coincidem quase todos
(\cod{PAY\_0}$\le1{,}5$, \cod{PAY\_2}$\le1{,}5$, \cod{PAY\_AMT3}$\le664{,}5$, \dots), como era
de esperar: as duas medidas de impureza são muito parecidas. A poda por \emph{log-loss} deixa 14 folhas (AUC $0{,}751$, contra
$0{,}611$ da cheia). O primeiro corte separa quem está com dois ou mais meses de atraso
em setembro --- o grupo de maior risco.

\textbf{(c)} Com 200 árvores: $m=\sqrt p$, AUC $0{,}773$ e Brier $0{,}1356$;
$m\approx p/3$, $0{,}772$; \emph{bagging}, $0{,}767$; folhas cheias, $0{,}760$ e Brier
$0{,}1387$ --- folhas de uma observação dão probabilidades mais extremas. O AdaBoost ordena
bem (AUC $0{,}771$), mas as probabilidades dele ficam entre $0{,}22$ e $0{,}67$: no corte
$1/6$ ele contata todo mundo, com custo $0{,}779$ --- o mesmo da regra trivial (o corte
por custo pode não cortar nada).
\end{solucao}

\begin{criterio}
(a) 0,15; (b) 0,2; (c) 0,25. O AdaBoost é opcional e não soma pontos, mas uma boa
leitura dele pode compensar até 0,1 perdido nesta tarefa.
\end{criterio}

%------------------------------------------------------------------------------
\tarefa{C7}{Escolha final e teste}{0,8 ponto}
%------------------------------------------------------------------------------

\begin{itens}
  \item Monte uma tabela, com probabilidades fora da dobra no treino, de todos os
        modelos: AUC, precisão média, Brier, custo no corte $1/6$, melhor corte e
        custo nele. Escolha o modelo final \emph{e} o corte, e justifique. Qual das
        colunas responde à pergunta do banco?
  \item Reajuste o modelo escolhido em todo o treino e aplique-o \textbf{uma vez} ao
        teste, no corte escolhido: a matriz de confusão, o custo médio por cliente, a
        sensibilidade e a fração de clientes contatados. Compare com as regras
        triviais.
  \item Faça um \emph{bootstrap} \textbf{pareado} no teste ($B=1000$,
        \cod{np.random.default\_rng(2026)}): em cada réplica, sorteie clientes com
        reposição e calcule, nos \emph{mesmos} clientes, a diferença de AUC e a
        diferença de custo entre o seu modelo e a logística com \cod{PAY\_*} em
        \emph{dummies}. Dê os intervalos de 95\% (percentis). O seu modelo é
        melhor que a logística? Por que o \emph{bootstrap} precisa ser pareado?
  \item Em um parágrafo, para a gerente de risco: quantos clientes contatar, quantos
        inadimplentes isso pega, e o que se ganha em relação a contatar todos.
\end{itens}

\begin{solucao}
\textbf{(a)} Pelo custo no corte $1/6$: floresta $m=\sqrt p$ $0{,}565$, floresta
$p/3$ $0{,}571$, \emph{bagging} $0{,}579$, logística com \emph{dummies} $0{,}580$, árvore
podada $0{,}586$, KNN $0{,}598$ --- todos bem abaixo de $0{,}779$. A coluna que
responde ao banco é o custo; a AUC ordena modelos sem saber dos custos, e o Brier diz
se o corte teórico é confiável. A floresta $m=\sqrt p$ tem o menor custo, e o melhor
corte dela ($0{,}16$) coincide com o teórico --- o que autoriza usar $1/6$. A logística
com \emph{dummies}, a $0{,}015$ de distância e muito mais interpretável, é uma escolha
igualmente defensável.

\textbf{(b)} Floresta no teste, corte $1/6$: VN 3632, FP 2209, FN 386, VP 1273. Custo
médio $0{,}552$ por cliente, contra $0{,}779$ de contatar todos (29\% menos);
contata $46{,}4\%$ dos clientes e pega $76{,}7\%$ dos inadimplentes. AUC $0{,}787$,
Brier $0{,}134$.

\textbf{(c)} AUC: floresta $0{,}787$, logística $0{,}784$, diferença com IC
$(-0{,}006;\ 0{,}012)$. Custo: floresta $0{,}552$, logística $0{,}538$, diferença com IC
$(-0{,}008;\ 0{,}035)$. Os dois intervalos contêm zero: o teste \emph{não distingue} os
dois modelos, e no custo a logística até ficou à frente nesta amostra. O
\emph{bootstrap} é pareado porque os dois modelos são avaliados nos mesmos clientes:
os erros são correlacionados (clientes difíceis são difíceis para os dois), e a
diferença tem variância muito menor do que a soma das variâncias. Reamostrar cada modelo
separadamente jogaria fora essa correlação e produziria intervalos largos demais.

\textbf{(d)} ``Contatando os clientes com probabilidade estimada de inadimplência acima
de 1 em 6 (cerca de 17\%) --- uns 46\% da carteira ---, pegamos três de cada quatro que vão ficar inadimplentes, e o
custo total cai cerca de 30\% em relação a contatar todos. Um modelo mais simples, a
regressão logística, faz praticamente o mesmo.''
\end{solucao}

\begin{criterio}
(a) 0,25 --- a escolha do corte e a justificativa pelo custo; (b) 0,2; (c) 0,25 ---
0,1 pelo cálculo pareado, 0,15 pela conclusão e pela explicação do pareamento;
(d) 0,1. Não há modelo ``certo'': a nota é da coerência entre a tabela, a escolha e a
conclusão.
\end{criterio}

%==============================================================================
\section*{Para fechar}
%==============================================================================
Se você chegou até aqui, percorreu o curso inteiro --- risco, viés e variância,
regularização, validação cruzada, métodos não paramétricos, árvores e florestas,
\emph{pipelines}, os classificadores generativos e discriminativos, as métricas, a SVM
--- em dados que ninguém inventou. Repare no que decidiu os resultados mais
importantes: quase nunca foi a escolha do método. Foi olhar os dados antes (as
duplicatas, o traço, os códigos do \cod{PAY}), avaliar do jeito que responde à
pergunta (por grupo, com o custo certo) e desconfiar dos números que pareciam bons
demais. É isso que se leva deste curso.

\end{document}
